%------------------------------------------------------------------------------%
\documentclass[a4paper,12pt,fleqn]{article}

\usepackage{integracao_series}

\ShowAnswers

%------------------------------------------------------------------------------%
\begin{document}

\makeHeader{Integração e Séries}{Prova 2 -- Turma 3}{27/04/2026}

~\\[2mm]
Calcule as integrais solicitadas, apresentando todos os passos necessários.

% 01
%------------------------------------------------------------------------------%
\exercise{25\%}
\(
  \int \sqrt{9 - x^2}\, dx
\)

\begin{answer}
\begin{multicols}{2}
  Utilizamos a substituição trigonométrica
  \[
    x = 3\sin\theta
    \qquad \qquad
    dx = 3\cos\theta\, d\theta
  \]
  temos
  \[
    \sqrt{9 - x^2}
    = \sqrt{9 - 9\sin^2\theta}
    = 3\cos\theta
  \]
  Substituindo na integral
  \begin{align*}
    F
    & = \int \sqrt{9 - x^2}\, dx                        \\[2mm]
    & = \int (3\cos\theta)(3\cos\theta)\, d\theta       \\[2mm]
    & = 9 \int \cos^2\theta\, d\theta
  \end{align*}
  Usamos a identidade
  \[
    \cos^2\theta = \frac{1 + \cos(2\theta)}{2}
  \]
  temos
  \begin{align*}
    F
    & = \frac{9}{2} \int 1 + \cos 2\theta\, d\theta   \\[2mm]
    & = \frac{9}{2} \left(\theta + \frac{1}{2}\sin 2\theta \right) + C
  \end{align*}
  Voltando para $x$ temos
  \begin{align*}
    \theta
    & = \arcsin\left(\frac{x}{3}\right)
    \\
    \sin 2\theta
    & = 2\sin\theta\cos\theta                        \\[2mm]
    & = 2\,\frac{x}{3}\,\frac{\sqrt{9 - x^2}}{3}     \\[2mm]
    & = \frac{2}{9}x\sqrt{9 - x^2}
  \end{align*}
  Assim
  \begin{align*}
    F(x)
    & = \frac{9}{2}
      \left(
        \arcsin\left(\frac{x}{3}\right)
        + \frac{1}{2}\frac{2}{9}x\sqrt{9 - x^2}
      \right) + C \\[2mm]
    & = \frac{9}{2}\arcsin\left(\frac{x}{3}\right) + \frac{x}{2}\sqrt{9 - x^2} + C
  \end{align*}
\end{multicols}
\end{answer}

% 02
%------------------------------------------------------------------------------%
\exercise{25\%}
\(
  \int_0^{\infty} \frac{e^x}{(1 + e^x)^2} \, dx
\)

\begin{answer}
\begin{multicols}{2}
Trata-se de uma integral imprópria com domínio infinito,
portanto precisamos explicitar o limite
\begin{align*}
  A
  & = \int_0^{\infty} \frac{e^x}{(1 + e^x)^2}\, dx \\[2mm]
  & = \lim_{b \to \infty} \int_0^b \frac{e^x}{(1 + e^x)^2}\, dx
\end{align*}
Primeiro vamos calcular a primitiva
\[
  F(x) = \int \frac{e^x}{(1 + e^x)^2}\, dx
\]
usando a substituição
\[
  u = 1 + e^x
  \qquad\qquad
  du = e^x dx
\]
temos
\begin{align*}
  F(x)
  & = \int \frac{1}{u^2}\, du   \\[2mm]
  & = \int u^{-2}\, du          \\[2mm]
  & = \frac{u^{-1}}{-1} + C     \\[2mm]
  & = \frac{-1}{u} + C     \\[2mm]
  & = \frac{-1}{1 + e^x} + C
\end{align*}
Calculamos agora a integral definida no intervalo~$[0, b]$
\begin{align*}
  I(b)
  & = \int_0^b \frac{e^x}{(1 + e^x)^2}\, dx       \\[2mm]
  & = F(x) \evalat{0}{b}                          \\[2mm]
  & = \frac{-1}{1 + e^x}\, \evalat{0}{b}          \\[2mm]
  & = \frac{-1}{1 + e^b} - \frac{-1}{1 + e^0}     \\[2mm]
  & = \frac{-1}{1 + e^b} + \frac{1}{1 + 1}        \\[2mm]
  & = \frac{1}{2} + \frac{1}{1 + e^b}
\end{align*}
Agora  calculamos o limite
\begin{align*}
  A
  & = \lim_{b \to \infty} I(b) \\[2mm]
  & = \lim_{b \to \infty} \left(\frac{1}{2} + \frac{1}{1 + e^b}\right) \\[2mm]
  & = \frac{1}{2}
\end{align*}
\end{multicols}
\end{answer}

% 03
%------------------------------------------------------------------------------%
\exercise{25\%}
\(
  \int_{-1}^1 x e^x \, dx
\)

\begin{answer}
\begin{multicols}{2}
  Queremos calcular
  \[
    A = \int_{-1}^1 x e^x \, dx
  \]
  Usamos integração por partes e escolhemos
  \[
    u = x
    \qquad\qquad
    dv = e^x dx
  \]
  obtermos assim
  \[
    du = dx
    \qquad\qquad
    v = e^x
  \]
  Calculando a primitiva
  \begin{align*}
    F(x)
    & = \int x e^x dx           \\*[2mm]
    & = x e^x - \int e^x dx     \\*[2mm]
    & = x e^x - e^x + C
  \end{align*}
  Calculando a integral definida
\begin{align*}
    A
    & = \int_0^1 x e^x dx                        \\[2mm]
    & = F(x) \evalat{-1}{1}                      \\[2mm]
    & = \left(x e^x - e^x\right) \evalat{-1}{1}  \\[2mm]
    & = \left(1 e^1 - e^1\right)
      - \left(-1 e^{-1} - e^{-1}\right)          \\[2mm]
    & = e - e + \frac{1}{e} + \frac{1}{e}        \\[2mm]
    & = \frac{2}{e}
\end{align*}
\end{multicols}
\end{answer}

% 04
%------------------------------------------------------------------------------%
\exercise{25\%}
  \(
    \int \sin^3 x \cos^2 x \, dx
  \)

\begin{answer}
\begin{multicols}{2}
  Queremos calcular a integral indefinida
  \[
    F(x)
    = \int \sin^3 x \cos^2 x \, dx
  \]
  Como temos uma potência ímpar no seno,\\
  realizamos a transformação
  \begin{align*}
    \sin^3 x
    & = \sin^2 x \sin x           \\
    & = \left(1 - \cos^2 x \right)\sin x
  \end{align*}
  Então
  \begin{align*}
    F(x)
    & = \int \sin^3 x \cos^2 x dx                      \\[2mm]
    & = \int \left(1 - \cos^2 x\right)\cos^2 x \sin x\, dx
  \end{align*}
  Usando a substituição
  \[
    u = \cos x
    \qquad\qquad
    du = -\sin x dx
  \]
  temos
  \begin{align*}
    F(x)
    & = \int \left(1 - u^2\right) u^2 (-du)          \\[2mm]
    & = \int u^4 - u^2 du                            \\[2mm]
    & = \frac{u^5}{5} - \frac{u^3}{3}  + C           \\[2mm]
    & = -\frac{\cos^3 x}{3} + \frac{\cos^5 x}{5} + C \\[2mm]
    & =  \frac{\cos^5 x}{5} - \frac{\cos^3 x}{3} + C
  \end{align*}
\end{multicols}
\end{answer}

%  \section*{Tabela de Integrais}
%
%  \[
%  \begin{array}{|c|c|}
%    \hline
%    \text{Função} & \text{Integral} \\
%    \hline
%    \int x^n dx & \frac{x^{n+1}}{n+1} + C \quad (n \neq -1) \\
%    \hline
%    \int \frac{1}{x} dx & \ln|x| + C \\
%    \hline
%    \int e^x dx & e^x + C \\
%    \hline
%    \int \sin x dx & -\cos x + C \\
%    \hline
%    \int \cos x dx & \sin x + C \\
%    \hline
%    \int \sec^2 x dx & \tan x + C \\
%    \hline
%    \int \frac{1}{\sqrt{a^2 - x^2}} dx & \arcsin\left(\frac{x}{a}\right) + C \\
%    \hline
%    \int \frac{1}{a^2 + x^2} dx & \frac{1}{a}\arctan\left(\frac{x}{a}\right) + C \\
%    \hline
%    \int \cos^2 x dx & \frac{x}{2} + \frac{\sin 2x}{4} + C \\
%    \hline
%    \int \sin^2 x dx & \frac{x}{2} - \frac{\sin 2x}{4} + C \\
%    \hline
%  \end{array}
%  \]

%------------------------------------------------------------------------------%
\end{document}
%------------------------------------------------------------------------------%
